\section{Related Work}

\textbf{Self-supervised Learning.} Modern SSL generally falls into two paradigms. \textit{Joint Embedding} (JE) methods, such as the MoCo~\citep{he2020momentum} and DINO~\citep{oquab2023dinov2} families, prioritize global invariance via multi-view alignment, yielding strong semantics but often losing spatial grounding. Conversely, \textit{Masked Image Modeling} (MIM), exemplified by MAE~\citep{he2022masked} and SimMIM~\citep{xie2022simmim}, emphasizes spatial equivariance through pixel reconstruction. While hybrids like iBOT~\citep{zhou2021ibot} and DINOv2~\citep{oquab2023dinov2} attempt to combine these objectives, they still rely heavily on global invariance and forgo pixel reconstruction. 

\textbf{Sparse Representation.} A growing body of work replaces dense feature maps with compact embeddings. Sparse R-CNN~\citep{sun2021sparse} and Mask2Former~\citep{cheng2022masked} utilize sparse queries for supervised tasks, while BLIP-2~\citep{li2023blip} and TiTok~\citep{yu2024image} employ sparse tokens for vision--language or generative efficiency. SemMAE~\citep{li2022semmae} utilizes sparse tokens to guide masking using a pretrained teacher. Unlike these methods, STELLAR treats sparse tokens as the \textbf{primary latent representation} and learns in SSL manner.

\textbf{Disentanglement \& Low-rank Factorization.} The assumption that high-dimensional data lie on low-dimensional manifolds is foundational to dictionary learning~\citep{mairal2008supervised}. In deep learning, low-rank constraints are typically applied to weights for efficiency (e.g., LoRA~\citep{hu2022lora}). STELLAR differs by applying \textit{low-rank factorization to the feature map itself}, disentangling "what" (semantic latents $\mL$) from "where" (spatial assignments $\mS$).



\paragraph{The Empirical Dilemma.} Current vision frameworks face a persistent gap: models excelling at pixel-level reconstruction often produce weaker semantic representations~\citep{zhang2022improving, chen2024deconstructing}, while those achieving top-tier semantics often abandon reconstruction to avoid low-level shortcuts~\citep{assran2023self, darcet2025cluster}. We demonstrate that by factorizing the latent representation, it is possible to achieve strong performance on both image understanding and reconstruction.


% \section{Related Work}
% \textbf{Self-supervised Learning} Recent progress in self-supervised visual learning largely focuses on relating global and local views of an image. Contrastive methods such as SimCLR~\citep{chen2020simple} and the MoCo series~\citep{he2020momentum, chen2020improved, chen2021empirical} highlighted the importance of large-scale augmentation and negative sampling, with MoCo introducing a momentum encoder and queue to relax batch-size requirements for efficient training. SwAV~\citep{caron2020unsupervised} combined contrastive learning with online clustering, while BYOL~\citep{grill2020bootstrap} and SimSiam~\citep{chen2021exploring} showed that teacher–student or stop-gradient mechanisms suffice without negatives. The DINO family~\citep{caron2021emerging, oquab2023dinov2} further revealed emergent segmentation ability by aligning global and cropped views.

% In parallel, masked image modeling (MIM) emerged, inspired by masked language modeling. BEiT~\citep{bao2021beit} reconstructed discrete tokens from an external tokenizer, while MAE~\citep{he2022masked} and SimMIM~\citep{xie2022simmim} simplified the design for pixel reconstruction with asymmetric encoder–decoder. Many followup methods discarded pixel reconstruction and instead performing alignment in the latent space~\citep{zhou2021ibot, baevski2022data2vec, assran2022masked, dong2022bootstrapped, chen2022sdae, tao2023siamese, yang2024one, liu2024exploring}. Our work differs by exploring \textit{sparse} self-supervised representations, aiming to combine the semantic richness of contrastive/distillation methods with the reconstruction fidelity of MIM, while avoiding their reliance on dense feature grids.

% % Empirically, contrastive/distillation methods excel at global invariance and classification, whereas MIM objectives yield stronger dense features and robustness to distribution shift.



% \textbf{Sparse Representation} A growing body of work replaces dense feature maps with a compact set of embeddings for downstream tasks. Sparse R-CNN~\citep{sun2021sparse}, DETR~\citep{carion2020end}, MaskFormer~\citep{cheng2021per} and Mask2Former~\citep{cheng2022masked} learn sparse queries from detection or segmentation supervision. BLIP-2~\citep{li2023blip} introduces a Q-Former to extract sparse tokens for efficient vision–language modeling, trained with paired image–text data. More recently, TiTok~\citep{yu2024image} represents images with sparse 1D tokens for efficient reconstruction and generation. TokenLearner~\citep{ryoo2021tokenlearner} integrate the sparse token idea into the ViT architecture, improving model efficiency by reducing the number of tokens. SemMAE~\citep{li2022semmae} uses a pretrainined iBOT model and learns part tokens to guided the masking process in MAE. In contrast, we treat the sparse tokens as the latent vision representation, and trained them for both low level reconstruction and high-level semantics in a fully self-supervised manner, without requiring annotations or image–text supervision.


% As a related concept, sparsity traditionally refers to vectors or matrices with mostly zero entries, and has been extensively studied in signal processing and machine learning. Classical approaches include compressed sensing~\citep{donoho2006compressed}, sparse coding and dictionary learning~\citep{olshausen1997sparse, mairal2008supervised, tovsic2011dictionary}, and regularization techniques such as the Lasso~\citep{tibshirani1996regression}. Sparsity has also been exploited in neural networks through sparse autoencoders~\citep{ng2011sparse} and structured pruning methods~\citep{anwar2017structured}. Our notion of sparse representation differs: rather than enforcing zeros in dense feature maps, we learn a compact set of informative tokens, each encoding semantically meaningful features for downstream tasks.

% \textbf{Low-rank Representation} Low-rank representation is a long-standing idea in machine learning, typically assuming that high-dimensional data lie on a low-dimensional manifold. Classical approaches include PCA, low-rank matrix recovery~\citep{candes2009accurate}, and dictionary learning or sparse coding~\citep{olshausen1997sparse, mairal2008supervised, tovsic2011dictionary}, where signals are expressed through a small set of basis vectors. In deep learning, low-rank constraints have been widely applied for efficiency: for example, low-rank factorization of neural network weights~\citep{sainath2013low}, low-rank approximations of attention maps~\citep{katharopoulos2020transformers}, and parameter-efficient fine-tuning methods such as LoRA~\citep{hu2022lora}. In contrast to these works, STELLAR applies low-rank factorization directly to the feature map from a single image, disentangling spatial and semantic information.

% This reveals an empirical dilemma in current vision frameworks: models that excel at pixel-level reconstruction often produce weaker semantic representations~\citep{zhang2022improving, chen2024deconstructing}. Conversely, state-of-the-art SSL methods that achieve strong semantics typically abandon pixel reconstruction to avoid low-level shortcuts~\citep{zhou2021ibot, assran2023self, darcet2025cluster}. In contrast, we demonstrate that by disentangling semantic and localization information, it is possible to learn sparse representations that simultaneously achieve strong image understanding and high-quality reconstruction.
